\section{Spinful Species, Raw Currents, and the Parity Sublattice}
\label{sec:ch21-spin-charge}

\subsection*{From two chiral branches to four fermion species}

For spinless fermions, the two chiral branches carry separate density
fluctuations.  Adding an electron spin label produces four branches, but it
also raises a question which is invisible if one starts immediately with
continuum ``charge'' and ``spin'' fields: \emph{which pairs of charge and spin
quantum numbers actually describe finite fermionic states?}  A second question
concerns the familiar claim of spin--charge separation.  The operator algebra
will separate into simple linear combinations, whereas commutation on a
finite energy budget and factorization of the state space require distinct
arguments.  We first formulate these questions entirely in the finite CAR
(carrier) representation, before using any oscillator interpretation.

Fix an integer $h\geq 1$, and write
\begin{equation}
  L=2h,\qquad
  \mathcal B_L=(-h,h]\cap\mathbb Z=\{1-h,\ldots,h\},\qquad
  \Lambda_L=\mathbb Z/L\mathbb Z.
  \label{eq:ch21-finite-band}
\end{equation}
Thus $L$ is the number of lattice sites and also the number of allowed
momenta per species.  We reserve $M\geq 1$ for a \emph{positive-mode cutoff}:
when a collection of modes is considered, $1\leq m\leq M$; an individual
signed density shift $m$ may instead be any integer.  The Fourier root and
orientation are fixed by
\begin{equation}
  \zeta=\exp(2\pi i/L),\qquad
  c_{\alpha,x}=\frac{1}{\sqrt L}
     \sum_{k\in\mathcal B_L}\zeta^{kx}c_{\alpha,k},
  \qquad x\in\Lambda_L.
  \label{eq:ch21-fourier-orientation}
\end{equation}
We use periodic boundary conditions.  The inverse transform carries
$\zeta^{-kx}$.  The finite \emph{non-wrapping} density modes below must not be
silently identified with the complete spatial Fourier coefficients of
$c_{\alpha,x}^{\dagger}c_{\alpha,x}$: the latter also contain momenta
wrapped across the ends of $\mathcal B_L$ (see Chapter~\ref{ch:density_modes}).

\begin{definition}[Ordered spinful CAR carrier]
\label{def:ch21-species}
Let
\begin{equation}
  \mathcal C=\{L,R\}\times\{\uparrow,\downarrow\}.
  \label{eq:ch21-species}
\end{equation}
The chiral label is $\eta\in\{L,R\}$ and the spin label is
$\sigma\in\{\uparrow,\downarrow\}$.  For explicit sign computations we
choose once and for all the total order
\begin{equation}
  (L,\uparrow)<(L,\downarrow)<(R,\uparrow)<(R,\downarrow),
  \label{eq:ch21-species-order}
\end{equation}
followed by increasing momentum within each species.  The one-particle mode
set is $\mathcal I=\mathcal C\times\mathcal B_L$.  The fermionic Hilbert
space has the orthonormal occupation basis
\begin{equation}
  \mathcal H_L=\ell^2\!\bigl(\mathcal P(\mathcal I)\bigr),
  \quad |S\rangle,\ S\subseteq\mathcal I,
  \qquad \dim\mathcal H_L=2^{4L}.
\end{equation}
For $i\in\mathcal I$ write $c_i$ and $c_i^\dagger$ for the
annihilation and creation maps.  With $\{A,B\}=AB+BA$ they obey
\begin{equation}
 \{c_i,c_j\}=\{c_i^\dagger,c_j^\dagger\}=0,
 \qquad \{c_i,c_j^\dagger\}=\delta_{ij}I.
 \label{eq:ch21-four-species-car}
\end{equation}
\end{definition}

The order in \eqref{eq:ch21-species-order} is not a spin anisotropy.  It is a
convention needed to assign signs to occupation kets, CAR words and, later,
Klein operators.  To see the mechanism, put
$\epsilon(i,S)=(-1)^{\#\{j\in S:j<i\}}$.  Then
\begin{align}
 c_i|S\rangle
   &=\begin{cases}\epsilon(i,S)|S\setminus\{i\}\rangle,&i\in S,\\
                     0,&i\notin S,\end{cases}\label{eq:ch21-ann-action}\\
 c_i^\dagger|S\rangle
   &=\begin{cases}\epsilon(i,S)|S\cup\{i\}\rangle,&i\notin S,\\
                     0,&i\in S.\end{cases}\label{eq:ch21-cre-action}
\end{align}
These formulas reproduce the CAR: for $i\ne j$ the two possible orders
of inserting or deleting distinct occupied modes have opposite signs,
whereas for $i=j$ exactly one of $c_ic_i^\dagger$ and
$c_i^\dagger c_i$ acts as the identity on each basis ket.
The construction and its general CAR proof were established in
Chapter~\ref{ch:car_fock}; its multispecies version and the canonical
occupation tensor coordinates are supplied by the Core A05 module.
This matters because fermion operators of different species
\emph{anticommute}, even though the corresponding even density operators
will commute.

Let $n_{\alpha,k}=c_{\alpha,k}^\dagger c_{\alpha,k}$ and define the
sea-subtracted species-charge observable
\begin{equation}
  \widehat N_\alpha=\sum_{k\in\mathcal B_L}n_{\alpha,k}-hI,
  \qquad \alpha\in\mathcal C.
  \label{eq:ch21-relative-species-charge}
\end{equation}
Here the reference sea occupies the $h$ momenta $k\leq 0$ in every
species.  Its relative charges vanish; an occupation configuration has
$N_\alpha\in[-h,h]\cap\mathbb Z$.  This subtraction is essential when
talking about zero modes: $\sum_k n_{\alpha,k}$ counts \emph{absolute}
particles, while $\widehat N_\alpha$ counts the integer charge relative
to the chosen sea.

For an admissible vector of species charges
$\mathbf N=(N_\alpha)_{\alpha\in\mathcal C}$ and an integer $K\geq0$, let
\begin{equation}
 \mathcal B_{\mathbf N,K}
  =\operatorname{span}\Bigl\{|S\rangle:
     N_\alpha(S)=N_\alpha\ \text{for all }\alpha,
     \ \mathcal E(S)=\sum_{\alpha\in\mathcal C}\mathcal E_\alpha(S)
          \leq K\Bigr\}.
 \label{eq:ch21-joint-budget}
\end{equation}
Each $\mathcal E_\alpha$ is the nonnegative, integer excitation grading
\emph{above the ground configuration of its own charge sector}, constructed
in Chapter~\ref{ch:vacuum_budget}.  Equation
\eqref{eq:ch21-joint-budget} defines a coordinate subspace, not an
operator-order inequality of the form $H\psi\leq K\psi$.
We write $P_{\mathbf N,K}$ for its orthogonal projection.  The budget
constrains the \emph{sum} of the four excitation energies; it does not
place a separate independent cutoff $K$ on every species.

\subsection*{Raw currents and their exact inverse}

For $m\in\mathbb Z$ and $\alpha\in\mathcal C$, define the ordered
non-wrapping density operator
\begin{equation}
  \rho_{\alpha,m}
   =\sum_{\substack{k\in\mathcal B_L\\k+m\in\mathcal B_L}}
        c_{\alpha,k+m}^{\dagger}c_{\alpha,k}.
  \label{eq:ch21-species-density}
\end{equation}
The sign convention is deliberate: $m>0$ raises the momentum and the
excitation grading, and
$\rho_{\alpha,m}^{\dagger}=\rho_{\alpha,-m}$.  Since
$\rho_{\alpha,m}$ is even in fermion operators, the global CAR imply
\begin{equation}
 [\rho_{\alpha,m},\rho_{\beta,n}]=0\quad(\alpha\ne\beta).
 \label{eq:ch21-cross-species-density}
\end{equation}
For example, moving a two-fermion word for $\rho_{\alpha,m}$ past
one for $\rho_{\beta,n}$ entails four minus signs.  No energy budget is
needed for \eqref{eq:ch21-cross-species-density}.

\begin{definition}[Unnormalized charge and spin currents]
\label{def:ch21-raw-currents}
At each chirality $\eta$ put
\begin{align}
 R^c_{\eta,m}&=\rho_{(\eta,\uparrow),m}
                 +\rho_{(\eta,\downarrow),m},
     \label{eq:ch21-raw-charge}\\
 R^s_{\eta,m}&=\rho_{(\eta,\uparrow),m}
                 -\rho_{(\eta,\downarrow),m}.
     \label{eq:ch21-raw-spin}
\end{align}
Here $R^{c,s}$ denotes a \emph{raw current}, not the right-moving
chirality $R$.  Neither current contains a factor $1/\sqrt2$.
\end{definition}

The change of variables is algebraically invertible over $\mathbb C$:
\begin{equation}
 \begin{pmatrix}R^c_{\eta,m}\\R^s_{\eta,m}\end{pmatrix}
 =\begin{pmatrix}1&1\\1&-1\end{pmatrix}
   \begin{pmatrix}\rho_{(\eta,\uparrow),m}\\
                   \rho_{(\eta,\downarrow),m}\end{pmatrix},
 \qquad
 \begin{pmatrix}\rho_{(\eta,\uparrow),m}\\
                   \rho_{(\eta,\downarrow),m}\end{pmatrix}
 =\frac12\begin{pmatrix}1&1\\1&-1\end{pmatrix}
   \begin{pmatrix}R^c_{\eta,m}\\R^s_{\eta,m}\end{pmatrix}.
 \label{eq:ch21-raw-inverse}
\end{equation}
Indeed, the square of the displayed integer matrix is $2I$.
Taking adjoints and using that its entries are real proves
$(R^{c,s}_{\eta,m})^\dagger=R^{c,s}_{\eta,-m}$.  The factors
$1/2$ in the inverse are harmless for operators over $\mathbb C$;
their meaning for \emph{integer charge eigenvalues} will be quite
different.

\subsection*{The edge obstruction and the finite common-margin theorem}

A tempting argument says that up- and down-spin densities commute and
therefore so do the charge and spin currents.  The conclusion does
\emph{not} follow.  The two same-spin commutators still remain, and a
finite band has edge-dependent corrections to the usual current algebra.
Define
\begin{equation}
 E_{\eta,\sigma}(m,n)
  =[\rho_{(\eta,\sigma),m},\rho_{(\eta,\sigma),n}].
 \label{eq:ch21-edge-operator}
\end{equation}
These are genuine finite-dimensional operators, not a priori scalar
multiples of the identity.

\begin{proposition}[Global mixed-current identity]
\label{prop:ch21-global-edge}
For every $m,n\in\mathbb Z$ and either chirality $\eta$,
\begin{align}
 [R^c_{\eta,m},R^s_{\eta,n}]
   &=E_{\eta,\uparrow}(m,n)-E_{\eta,\downarrow}(m,n),
       \label{eq:ch21-mixed-edge}\\
 [R^c_{\eta,m},R^c_{\eta,n}]
   &=[R^s_{\eta,m},R^s_{\eta,n}]
     =E_{\eta,\uparrow}(m,n)+E_{\eta,\downarrow}(m,n).
       \label{eq:ch21-self-edge}
\end{align}
If $\eta\ne\eta'$, every commutator between
$R^a_{\eta,m}$ and $R^b_{\eta',n}$ vanishes for
$a,b\in\{c,s\}$, as an identity on all of $\mathcal H_L$.
\end{proposition}
\begin{proof}
Use bilinearity of $[A,B]=AB-BA$.  For instance, suppressing
$\eta,m,n$ when unambiguous,
\begin{align*}
 [R^c_m,R^s_n]
  &=[\rho_{\uparrow,m}+\rho_{\downarrow,m},
      \rho_{\uparrow,n}-\rho_{\downarrow,n}]\\
  &=[\rho_{\uparrow,m},\rho_{\uparrow,n}]
    -[\rho_{\uparrow,m},\rho_{\downarrow,n}]
    +[\rho_{\downarrow,m},\rho_{\uparrow,n}]
    -[\rho_{\downarrow,m},\rho_{\downarrow,n}].
\end{align*}
The middle two terms vanish by
\eqref{eq:ch21-cross-species-density}, giving
\eqref{eq:ch21-mixed-edge}.  With two plus signs, or two minus signs,
the same expansion gives \eqref{eq:ch21-self-edge}.  For different
chiralities all four species in the expansion are distinct, so again
every constituent commutator vanishes.  No rearrangement of fermionic
\emph{odd} factors was used.
\end{proof}

The single-species Schwinger coefficient in the convention of this book is
\begin{equation}
 \kappa(m,n)=
 \begin{cases}-m,&m+n=0,\\0,&m+n\ne0.\end{cases}
 \label{eq:ch21-schwinger-coefficient}
\end{equation}
It is the coefficient of a commutator acting on an appropriate budget,
not a globally valid nonzero scalar CCR on the finite Fock space.  The
origin of this distinction is the exact diagonal edge formula and the
frozen-margin theorem of Chapters~\ref{ch:heisenberg_algebra} and
\ref{ch:vacuum_budget}; their multispecies current version is provided
in the Core A06 support module.

\begin{theorem}[Raw charge/spin current algebra on a common budget]
\label{thm:ch21-budget-ccr}
Assume $h>0$, $K\geq0$, and an admissible charge vector
$\mathbf N$.  Fix $\eta\in\{L,R\}$ and $m,n\in\mathbb Z$.
Suppose \emph{both} spin species at that chirality satisfy
\begin{equation}
 \begin{split}
 |m|+|n|+K+|N_{(\eta,\uparrow)}|&\leq h,\\
 |m|+|n|+K+|N_{(\eta,\downarrow)}|&\leq h.
 \end{split}
 \label{eq:ch21-common-margin}
\end{equation}
For every $v\in\mathcal B_{\mathbf N,K}$, with the commutators
formed in the \emph{ambient} finite CAR representation,
\begin{align}
 [R^c_{\eta,m},R^s_{\eta,n}]v&=0,
   \label{eq:ch21-cross-ccr-budget}\\
 [R^c_{\eta,m},R^c_{\eta,n}]v
   &=[R^s_{\eta,m},R^s_{\eta,n}]v
     =2\kappa(m,n)v.
   \label{eq:ch21-self-ccr-budget}
\end{align}
In particular, if $m\geq1$ and
$2m+K+\max_{\sigma}|N_{(\eta,\sigma)}|\leq h$, then
\begin{equation}
 [R^c_{\eta,-m},R^c_{\eta,m}]v
   =[R^s_{\eta,-m},R^s_{\eta,m}]v=2m\,v,
 \qquad [R^c_{\eta,-m},R^s_{\eta,m}]v=0.
 \label{eq:ch21-two-m}
\end{equation}
\end{theorem}
\begin{proof}
The established single-species signed-current theorem states that,
under the corresponding margin in \eqref{eq:ch21-common-margin},
\begin{equation*}
 E_{\eta,\sigma}(m,n)v=\kappa(m,n)v
 \qquad(\sigma=\uparrow,\downarrow).
\end{equation*}
The same vector $v$ belongs to the joint budget, so each of these
identities applies \emph{to that vector}; the two separate conditions
are needed because the relative species charges can differ.  Substitute
them into \eqref{eq:ch21-mixed-edge} to obtain
$(\kappa-\kappa)v=0$.  Substituting into
\eqref{eq:ch21-self-edge} gives $(\kappa+\kappa)v=2\kappa v$.
Finally, taking the ordered pair of signed shifts $(-m,m)$ gives
$\kappa(-m,m)=m$ and $|{-m}|+|m|=2m$.
This proves all statements without postulating a finite-dimensional
Heisenberg representation.
\end{proof}

For all $1\leq m,n\leq M$, the simple \emph{sufficient} uniform
condition
\begin{equation}
 2M+K+\max_{\substack{\eta\in\{L,R\}\\
                         \sigma\in\{\uparrow,\downarrow\}}}
                   |N_{(\eta,\sigma)}|\leq h
 \label{eq:ch21-uniform-margin}
\end{equation}
allows the theorem to be applied to every pair of shifts among
$\{-M,\ldots,M\}$ at both chiralities.  It is not claimed to be
necessary.  In particular, a proof involving a \emph{product of several
currents} must additionally account for the intermediate excursions;
checking a one-commutator bound does not certify a whole word.

\begin{example}[The first nontrivial charge/spin oscillator]
\label{ex:ch21-first-mode}
Take $h=2$, so $L=4$ and
$\mathcal B_L=\{-1,0,1,2\}$, with all species charges zero and
$K=0$.  The ground ket $\Omega$ occupies $-1$ and $0$ in each
species.  For a fixed chirality and spin, the only unblocked
$m=1$ hop from this ground configuration is $0\to1$; write
$|e_\sigma\rangle=\rho_{(\eta,\sigma),1}\Omega$.
The CAR sign formula gives
\begin{equation*}
 \rho_{(\eta,\sigma),-1}\Omega=0,
 \qquad \rho_{(\eta,\sigma),-1}|e_\sigma\rangle=\Omega,
 \qquad \langle e_\uparrow,e_\downarrow\rangle=0.
\end{equation*}
Thus $R^c_{\eta,1}\Omega=|e_\uparrow\rangle+|e_\downarrow\rangle$
whereas $R^s_{\eta,1}\Omega=|e_\uparrow\rangle-|e_\downarrow\rangle$.
Their squared norms are both $2$, and direct application gives
\begin{equation}
 [R^c_{\eta,-1},R^c_{\eta,1}]\Omega=2\Omega,
 \quad [R^s_{\eta,-1},R^s_{\eta,1}]\Omega=2\Omega,
 \quad [R^c_{\eta,-1},R^s_{\eta,1}]\Omega=0.
 \label{eq:ch21-example-two}
\end{equation}
The bounds hold because $2(1)+0+0\leq h$.  Notice that the
\emph{positive} shift creates an excitation in either spin species:
the relative plus or minus sign distinguishes charge and spin
combinations, while both retain the raw normalization $2m$.
\end{example}

The budget hypothesis is substantive.  For example, keep one chirality,
put its up-spin species into the \emph{completely filled} band, its
down-spin species into the sea ground, and leave the other two species
in their sea grounds.  Every nonzero-momentum density hop annihilates
the full up-spin band, whereas the down-spin sector obeys
$[\rho_{\downarrow,-1},\rho_{\downarrow,1}]\Omega_\downarrow
 =\Omega_\downarrow$ for $h\geq2$.  Consequently on the resulting
nonzero product ket $v$,
\begin{equation}
  [R^c_{\eta,-1},R^s_{\eta,1}]v=-v\ne0.
  \label{eq:ch21-off-margin-example}
\end{equation}
The up-spin relative charge is $h$, so the common bound fails.
This is a direct counterexample to unrestricted global spin/charge
commutation on the finite CAR carrier.

There is another important boundary effect even \emph{inside} the
regime of Theorem~\ref{thm:ch21-budget-ccr}.  Set
$P=P_{\mathbf N,K}$, $Q=I-P$, and let $A,B$ be any ambient maps.
Insertion of $I=P+Q$ gives the exact identity
\begin{equation}
 [PAP,PBP]
   =P[A,B]P-PAQBP+PBQAP.
 \label{eq:ch21-compression-residual}
\end{equation}
Here the left side contains \emph{compressed} operators, unlike the
ambient commutators of the theorem.  In the preceding $K=0$ example,
$P$ has rank one, $PR^{c}_{\eta,\pm1}P=0$ and therefore the
compressed commutator is zero, although
$P[R^c_{\eta,-1},R^c_{\eta,1}]P=2P$.  The difference is precisely
supplied by the excursions through $Q\mathcal H_L$.  This example
makes concrete why no finite-dimensional compression can realize a
nonzero scalar CCR as an identity on its entire range: taking the
trace of $[A,B]=\lambda I$ in finite dimension would give
$0=\lambda\dim\mathcal H$.

\subsection*{Integer zero modes and the parity sublattice}

The same linear transformation must now be applied to the integer
\emph{zero-mode} charges, not to oscillator operators.  For each
chirality, define the commuting self-adjoint observables
\begin{equation}
 \widehat Q_{c,\eta}=\widehat N_{(\eta,\uparrow)}
                         +\widehat N_{(\eta,\downarrow)},
 \qquad
 \widehat Q_{s,\eta}=\widehat N_{(\eta,\uparrow)}
                         -\widehat N_{(\eta,\downarrow)}.
 \label{eq:ch21-zero-modes}
\end{equation}
They are called charge and spin zero modes in \emph{raw integer units};
$Q_s$ is twice the $S^z$ polarization in the usual spin-$\tfrac12$
normalization.  On $\mathcal B_{\mathbf N,K}$ their eigenvalues are
\begin{equation}
 Q_{c,\eta}=N_{(\eta,\uparrow)}+N_{(\eta,\downarrow)},
 \qquad
 Q_{s,\eta}=N_{(\eta,\uparrow)}-N_{(\eta,\downarrow)}.
 \label{eq:ch21-zero-eigenvalues}
\end{equation}
This follows immediately from the diagonal species-charge observable
and its established action on a fixed sector.

\begin{proposition}[Parity sublattice and integral inverse]
\label{prop:ch21-parity-lattice}
Consider the map of integer charge lattices
\begin{equation}
 T:\mathbb Z^2\longrightarrow\mathbb Z^2,
 \qquad T(a,b)=(a+b,a-b).
 \label{eq:ch21-parity-map}
\end{equation}
Its image is exactly the index-two subgroup
\begin{equation}
 \Lambda_{c,s}
 =\{(q_c,q_s)\in\mathbb Z^2:q_c\equiv q_s\pmod2\}.
 \label{eq:ch21-parity-sublattice}
\end{equation}
On that subgroup $T$ is a bijection, with inverse
\begin{equation}
 T^{-1}(q_c,q_s)
 =\left(\frac{q_c+q_s}{2},\frac{q_c-q_s}{2}\right).
 \label{eq:ch21-parity-inverse}
\end{equation}
The inclusion $\Lambda_{c,s}\subsetneq\mathbb Z^2$ is proper;
for example, $(1,0)\notin\Lambda_{c,s}$.
\end{proposition}
\begin{proof}
For $a,b\in\mathbb Z$,
$T(a,b)_c-T(a,b)_s=(a+b)-(a-b)=2b$;
hence both coordinates have the same parity.  Conversely, assume
$q_c\equiv q_s\pmod2$.  Then $(q_c-q_s)/2$ is an integer, and
$q_c+q_s=(q_c-q_s)+2q_s$ is also even.  Therefore the two coordinates
in \eqref{eq:ch21-parity-inverse} are integers.  Applying $T$
to that pair yields
\begin{equation*}
  \left(\frac{q_c+q_s}{2}+\frac{q_c-q_s}{2},
        \frac{q_c+q_s}{2}-\frac{q_c-q_s}{2}\right)=(q_c,q_s).
\end{equation*}
Applying the proposed inverse to $T(a,b)$ returns $(a,b)$ in the
same way.  This proves injectivity, surjectivity onto the displayed
image, and the inverse formula.  Finally $1-0$ is odd, so $(1,0)$
is excluded.  Equivalently, the integer matrix for $T$ has determinant
$-2$; its image has index two and is not all of $\mathbb Z^2$.
\end{proof}

For the full four-species problem, the proposition applies
\emph{separately} at $\eta=L$ and $\eta=R$.  In other words, the
physical integer labels belong to
$\Lambda_{c,s}^{(L)}\times\Lambda_{c,s}^{(R)}$, with additional
finite-band charge bounds inherited from $-h\leq N_\alpha\leq h$.
A single parity statement for the \emph{sum} of the two chiralities
would be weaker and would not recover both independent species-charge
pairs.  The parity rule does not project out fermionic states: it is an
arithmetic description of the image of the existing species labels
under a non-unimodular change of coordinates.

\begin{example}[An allowed zero mode and a forbidden formal label]
Take $h=2$ and one chirality with
$N_\uparrow=2$ and $N_\downarrow=-1$.  The raw labels are
$(Q_c,Q_s)=(1,3)$.  They have equal parity, and
\eqref{eq:ch21-parity-inverse} recovers
$(N_\uparrow,N_\downarrow)=(2,-1)$.  By contrast, a hypothetical
$(Q_c,Q_s)=(1,0)$ would give
$(N_\uparrow,N_\downarrow)=(1/2,1/2)$; it is a perfectly reasonable
point of $\mathbb R^2$ but not a finite fermionic charge sector.
This is why replacing integer charges by two unconstrained
independent bosonic zero modes changes the state-label set.
\end{example}

\subsection*{What factorizes, and what does not}

Because the one-particle labels of different species are disjoint, the
occupation basis gives a unitary \emph{ordered} identification of the
unrestricted finite Hilbert space with the tensor product of the four
single-species occupation spaces.  For odd CAR operators the ordering
requires fermionic parity strings; for even density operators their
action can be taken on the corresponding tensor factor.  This
occupation-space factorization is a theorem of the finite CAR
construction, but it is \emph{not} a factorization into independent
charge and spin zero-mode lattices.

To make the energy restriction explicit, denote by
$\mathcal H_{\alpha,N_\alpha}^{(e)}$ the subspace of a single
species with relative charge $N_\alpha$ and excitation energy exactly
$e\in\mathbb N$.  Regrouping the occupation basis of
\eqref{eq:ch21-joint-budget} yields the canonical graded-vector-space
identification
\begin{equation}
 \mathcal B_{\mathbf N,K}\ \simeq\!
 \bigoplus_{\substack{(e_\alpha)_{\alpha\in\mathcal C}\in\mathbb N^4\\
                    \sum_\alpha e_\alpha\leq K}}
 \bigotimes_{\alpha\in\mathcal C}
       \mathcal H_{\alpha,N_\alpha}^{(e_\alpha)}.
 \label{eq:ch21-budget-energy-sum}
\end{equation}
The direct sum is orthogonal, and any zero-dimensional energy-grade
factor contributes nothing.  To prove the identity, map an occupation
ket to the ordered tensor product of its four restrictions to species.
This is a bijection of orthonormal bases; by definition of the total
excitation grading, the ket belongs to the budget precisely when the
four resulting grades have sum at most $K$.  Linearity proves the
claimed unitary identification.  No statement about charge/spin
\emph{oscillator} tensor factors is required for this proof.

In general,
\begin{equation}
  \mathcal B_{\mathbf N,K}
  \not\simeq_{\mathrm{natural}}
  \bigotimes_{\alpha\in\mathcal C}
      \Bigl(\bigoplus_{e=0}^{K}
             \mathcal H_{\alpha,N_\alpha}^{(e)}\Bigr),
 \label{eq:ch21-not-product-cutoff}
\end{equation}
where the notation emphasizes that the \emph{natural inclusion} of
the left-hand space into the right-hand tensor product is generally
strict; it is not a statement that abstract finite-dimensional vector
spaces could never have accidentally equal dimensions.  For example,
at $K=1$ a state with one energy-one excitation in each of two
species lies in the product of separate $K=1$ cuts, but has total
energy two and is absent from the joint budget.  In particular, a
putative $\mathcal B_{c,K}\otimes\mathcal B_{s,K}$ cannot be inferred
from a change of current variables, even before addressing the parity
sublattice.

\subsection*{Exchange-symmetric quadratic couplings}

We now isolate the precise sense in which spin and charge \emph{do}
separate algebraically.  Let $X_\uparrow,X_\downarrow$ and
$Y_\uparrow,Y_\downarrow$ be operators on a common vector space.
They may be density modes at different chiralities or momenta, and
we do \emph{not} assume that $X$ and $Y$ commute.  Introduce the raw
combinations $X_c=X_\uparrow+X_\downarrow$,
$X_s=X_\uparrow-X_\downarrow$, and similarly for $Y$.

\begin{proposition}[Ordered quadratic block identity]
\label{prop:ch21-quadratic-block}
For arbitrary $a,b\in\mathbb C$,
\begin{align}
 &a(X_\uparrow Y_\uparrow+X_\downarrow Y_\downarrow)
  +b(X_\uparrow Y_\downarrow+X_\downarrow Y_\uparrow)
 \notag\\
 &\hspace{1.5cm}=
  \frac{a+b}{2}\,X_cY_c+
  \frac{a-b}{2}\,X_sY_s.
 \label{eq:ch21-quadratic-block}
\end{align}
Thus a quadratic coupling matrix of the exchange-symmetric form
$\begin{psmallmatrix}a&b\\b&a\end{psmallmatrix}$ is block diagonal
in the raw charge/spin variables, with the coefficients displayed
above.  The statement preserves the product order throughout.
\end{proposition}
\begin{proof}
Expanding only by distributivity,
\begin{align*}
 X_cY_c&=X_\uparrow Y_\uparrow+X_\uparrow Y_\downarrow
          +X_\downarrow Y_\uparrow+X_\downarrow Y_\downarrow,\\
 X_sY_s&=X_\uparrow Y_\uparrow-X_\uparrow Y_\downarrow
          -X_\downarrow Y_\uparrow+X_\downarrow Y_\downarrow.
\end{align*}
The coefficients of the two same-spin terms become
$((a+b)+(a-b))/2=a$; those of the cross-spin terms become
$((a+b)-(a-b))/2=b$.  At no point do we commute an $X$ past a $Y$.
Equivalently, the change-of-basis matrix in
\eqref{eq:ch21-raw-inverse} diagonalizes
$\begin{psmallmatrix}a&b\\b&a\end{psmallmatrix}$ over
$\mathbb C$, but the explicit expansion shows that the matrix
argument remains valid for noncommutative operator entries.
\end{proof}

For a quadratic current Hamiltonian built as a sum of such ordered
bilinears, with an exchange-symmetric spin matrix \emph{for each}
choice of chirality and momentum labels, the proposition may be
applied term by term.  One obtains a charge-current part and a
spin-current part with no mixed bilinear.  Hermiticity of the
Hamiltonian requires the usual conjugate/reversed-index conditions
on its coefficients separately; spin exchange symmetry alone is not
an adjointness hypothesis.

The hypothesis is also necessary to inspect rather than assume.
For instance, a spin-dependent diagonal coupling with coefficients
$a_\uparrow\ne a_\downarrow$ becomes
\begin{equation}
 a_\uparrow X_\uparrow Y_\uparrow
 +a_\downarrow X_\downarrow Y_\downarrow
 =\frac{a_\uparrow+a_\downarrow}{4}(X_cY_c+X_sY_s)
  +\frac{a_\uparrow-a_\downarrow}{4}(X_cY_s+X_sY_c).
 \label{eq:ch21-asymmetry-mixes}
\end{equation}
The last two terms mix charge and spin.  Even when
\eqref{eq:ch21-quadratic-block} holds, it implies neither global
commutation of its finite-CAR charge and spin Hamiltonian parts nor a
unitary factorization of a \emph{compressed} Hamiltonian.  Global
current commutators retain the edge difference of
\eqref{eq:ch21-mixed-edge}; products require further excursion
bounds, and compressions retain the residual terms of
\eqref{eq:ch21-compression-residual}.  Nonquadratic interactions,
including the operators considered in later sections, are not
covered by this proposition.

\subsection*{Do the charge and spin modes have different velocities?}

A quadratic block decomposition permits distinct effective mode
velocities but does not \emph{force} them.  To state the condition
without presupposing a continuum limit, consider for each formal
charge/spin oscillator block $\nu\in\{c,s\}$ the stable two-chirality
pairing Hamiltonian
\begin{equation}
 H_{\nu,m}
 =m\Bigl[v_{1,\nu}
   (a_{R,\nu}^{\dagger}a_{R,\nu}
    +a_{L,\nu}^{\dagger}a_{L,\nu})
  +v_{2,\nu}(a_{R,\nu}^{\dagger}a_{L,\nu}^{\dagger}
                         +a_{R,\nu}a_{L,\nu})\Bigr],
 \quad m\geq1.
 \label{eq:ch21-pairing-hamiltonian}
\end{equation}
Here the $a$'s satisfy exact CCR on an \emph{abstract oscillator
carrier}; this is a model for the budget-stable current algebra, not
a global canonical representation by the finite-dimensional CAR
operators.  Assume real coefficients and
$v_{1,\nu}>|v_{2,\nu}|$.  Using the CCR gives
\begin{equation}
 \left[\begin{pmatrix}a_{R,\nu}\\a_{L,\nu}^{\dagger}\end{pmatrix},
                  H_{\nu,m}\right]
 =m\begin{pmatrix}
  v_{1,\nu}&v_{2,\nu}\\-v_{2,\nu}&-v_{1,\nu}
 \end{pmatrix}
 \begin{pmatrix}a_{R,\nu}\\a_{L,\nu}^{\dagger}\end{pmatrix}.
 \label{eq:ch21-bogoliubov-dynamic}
\end{equation}
The characteristic polynomial of this $2\times2$ matrix is
$\lambda^2-(v_{1,\nu}^2-v_{2,\nu}^2)$; therefore its positive
frequency, divided by $m$, is
\begin{equation}
 u_\nu=\sqrt{v_{1,\nu}^2-v_{2,\nu}^2}>0.
 \label{eq:ch21-block-speed}
\end{equation}
This is an exact calculation \emph{within the stated quadratic CCR
model}.  It follows that
\begin{equation}
 u_c\ne u_s
 \quad\Longleftrightarrow\quad
 v_{1,c}^2-v_{2,c}^2\ne v_{1,s}^2-v_{2,s}^2,
 \label{eq:ch21-speed-nondegeneracy}
\end{equation}
provided both radicands are strictly positive.  The equivalence
uses only injectivity of the positive square root.  Its reverse
implication would require re-examination without stability and
positivity.

If one parametrizes a block by real hyperbolic numbers
$c_\nu,s_\nu$ satisfying $c_\nu^2-s_\nu^2=1$, then one may use
\begin{equation}
 v_{1,\nu}=u_\nu(c_\nu^2+s_\nu^2),\qquad
 v_{2,\nu}=2u_\nu c_\nu s_\nu,\qquad
 g_\nu=(c_\nu-s_\nu)^2,\qquad
 g_\nu^{-1}=(c_\nu+s_\nu)^2.
 \label{eq:ch21-hyperbolic-parameters}
\end{equation}
The last identity follows by multiplication of the two squares and
$c_\nu^2-s_\nu^2=1$.  The symbol $g_\nu$ denotes a quadratic-model
Luttinger parameter, \emph{not} the real Umklapp coupling $g_U$
used later.  We also reserve $K$ exclusively for the finite energy
budget, not for the Luttinger parameter.

\begin{example}[Exchange symmetry without velocity separation]
\label{ex:ch21-equal-speeds}
Start with a spin-independent free linearized model of velocity
$v>0$.  Its spin-exchange-symmetric quadratic coupling has no
up/down mixing, $b=0$, and the same diagonal coefficient $a=v$
for both spins.  Equation~\eqref{eq:ch21-quadratic-block} gives
$\tfrac v2 X_cY_c+\tfrac v2 X_sY_s$; it \emph{does} have charge
and spin blocks.  After consistent canonical oscillator
normalization, both blocks have
\begin{equation*}
  v_{1,c}=v_{1,s}=v,\qquad v_{2,c}=v_{2,s}=0,
  \qquad u_c=u_s=v,
  \qquad c_c=c_s=1,\quad s_c=s_s=0.
\end{equation*}
Hence symmetry and even a quadratic charge/spin decomposition do not
imply unequal propagation speeds.  For contrast, the stable parameter
choices $(v_{1,c},v_{2,c})=(17,15)$ and
$(v_{1,s},v_{2,s})=(5,0)$ give $u_c=8$ and $u_s=5$.
The latter distinction comes from the \emph{parameters}, not merely
from the existence of the charge/spin basis.
\end{example}

\begin{remark}[Mathematical and formalization scope]
\label{rem:ch21-formalization-scope}
Three separate conclusions have been established in this section.
The four-species CAR and the global density identities concern the
\emph{entire finite Fock carrier}; the coefficient $2m$ and mixed
current cancellation concern the \emph{ambient action on specified
frozen-margin budget vectors}; and the dispersion computation
concerns a \emph{specified abstract quadratic CCR model}.  None
alone identifies its oscillator vacuum with a finite interacting
fermionic ground state, proves a thermodynamic limit, or provides a
spectral gap.  A charge/spin oscillator tensor-product description
would also have to respect the parity lattice and the total-energy
cutoff.  Such a construction is not inferred here.

The finite multispecies CAR and its support results are present in
the supplied Core modules, but the Chapter~21 spin/charge statements
as a combined Lean development are \emph{not yet implemented or
verified}.  The explicit conventions to reconcile at whole-book
integration are the momentum/Fourier orientation, sea subtraction for
zero modes, raw-versus-normalized current factors, ambient-versus-
compressed products, and the inherited sign order
\eqref{eq:ch21-species-order}.  For physical background on the
charge/spin description of one-dimensional fermions, see
\cite{giamarchi2004,gogolin1998,vondelft1998}; these references
are context rather than certificates for the finite margin theorem.
\end{remark}
